\subsection{流形结构的逆像、子流形}

5.8.1. 设 X 是一个拓扑空间，Y 是一个流形，f 是从 X 到 Y 的映射。考虑下列条件：

(QR)（分别为 (R)）对于每个 \(a \in X\)，存在 X 中 a 的一个开邻域 \(U\) 以及流形 Y 在 \(X\) 处的一个图 \((V, \varphi, E)\)，使得 \(Y\)\(f(a)\)\(f(U) \subset V\)，且 \(\varphi \circ f\) 将 U 同胚地映到 \(U\)\(\varphi(V)\) 与 E 的一个闭向量子空间 \(F\) 的交集上（分别地，映到与 E 的一个具有拓扑补空间的闭向量子空间 \(E\) 的交集上）。

若条件 (QR) 成立，则 X 上存在唯一的流形结构，使得三元组 (U, \(\varphi \circ f\) , F)（记号如 (QR)）是 X 的图。称其为由 \(f\) 从 Y 的流形结构得到的逆像结构。其拓扑就是 X 的拓扑。

X 上存在一个与给定拓扑相容且使 \(f\) 成为浸入的流形结构的充分必要条件是条件 (R) 成立。此结构随后是唯一的：它是 Y 的流形结构在 f 下的逆像。

5.8.2. 特别地，当对每个 \(a \in X\) 都存在 a 的一个开邻域 U，使得 \(a\)\(f|U\) 是从 \(U\) 到 Y 的某个开子集上的同胚时，上述条件 (R) 成立。在此情形下，得到的流形 \(Y\)\(X\) 在 \(Y\) 上是 étale 的（5.7.6）。

5.8.3. 现在设 X 是 Y 的一个拓扑子空间，f 是典范嵌入。若 f 满足第 5.8.1 号的条件 (R)（分别为 (QR)），则赋予 X 以 Y 的流形结构在 f 下的逆像结构后，称 X 为 Y 的子流形（分别为拟子流形）。子流形是拟子流形。

每个拟子流形都是局部闭的；每个开集都是子流形，且其结构就是第 5.2.3 号所定义的结构。

当 Y 局部有限维时，Y 的拓扑子空间 X 是 Y 的子流形，当且仅当对于每个 \(a \in X\)，存在 Y 中 a 的一个开邻域 U、Y 在 U 中的一个坐标系 \((\xi^{1}, \ldots, \xi^{m})\) 以及一个整数 \(k \leqslant m\)，使得 \(U \cap X\) 是 U 中使得当 \(\xi^{i}\) 时函数 \(1 \leqslant i \leqslant k\) 同时为零的点集。

设 \(f:X\to Y\) 是流形的一个态射。假设 \(f\)\(f\) 是浸入，且 f 诱导出从 X 到 \(f(X)\) 的同胚。则 \(f(X)\) 是 Y 的子流形，且 f 诱导出从 X 到 \(f\)\(f(X)\) 的同构；此时称 \(f\) 是 X 到 Y 的嵌入。

5.8.4. 设 Y 和 W 是两个流形，X 是 Y 的一个子流形，f 是从 X 到 W 的映射。f 是态射的充分必要条件是：对于每个 \(a \in X\)，存在 Y 中 a 的一个邻域 U 和一个态射 \(g : U \to W\)，使得 g 在 \(U \cap X\) 上与 f 重合。

5.8.5. 设 X 是流形 Y 的一个拟子流形，g 是从流形 Z 到 X 的映射。假设 K 的特征为 0，或 X 是 Y 的子流形。g 是从 Z 到 X 的态射的充分必要条件是它是从 Z 到 Y 的态射；换言之，X 的流形结构由 Y 的流形结构诱导（集合论，第~IV~章，§ 2，第 4 号）。

5.8.6. 若 X 是 Y 的子流形（分别为拟子流形），且 Y 是流形 Z 的子流形，则 X 是 Z 的子流形（分别为拟子流形）。

5.8.7. 设 \(X_{i}\)（i = 1, 2）是一个流形，\(Y_{i}\) 是 \(X_{i}\) 的子流形（分别为拟子流形）（i = 1, 2）。则 \(Y_{1} \times Y_{2}\) 是 \(X_{1} \times X_{2}\) 的子流形（分别为拟子流形）。

5.8.8. 设 X 是流形 Y 的一个拟子流形；令 x 是 X 的一点，并以 f 表示 X 到 Y 的典范嵌入。映射 \(T_{x}(f):T_{x}(X)\to T_{x}(Y)\) 是单射，从而可以把空间 \(T_{x}(X)\) 与 \(T_{x}(Y)\) 的一个闭向量子空间等同起来。拓扑向量空间商 \(T_{x}(Y)/T_{x}(X)\) 是一个 Banach 空间，称为 X 在 x 处（于 Y 中）的横截空间。它的维数（有限或 +∞）称为 X 在点 x 处于 Y 中的余维。

若进一步 X 是 Y 的子流形，则空间 \(T_{x}(X)\) 在 \(T_{x}(Y)\) 中具有拓扑补空间。

5.8.9. 设 \(f\) 是从流形 X 到流形 Y 的映射，\(\Gamma\) 是它的图。\(f\) 是态射的充分必要条件是满足下列两个条件：

\begin{enumerate}
\def\labelenumi{(\roman{enumi})}
\item
  \(\Gamma\) 是 \(X \times Y\) 的子流形。
\item
  对每个 \((x, y) \in \Gamma\)，有
\end{enumerate}

\[
\mathrm{T} _ {(x, y)} (\mathrm{X} \times \mathrm{Y}) = \mathrm{T} _ {(x, y)} (\Gamma) \oplus \mathrm{T} _ {y} (\mathrm{Y}).
\]

若此条件成立，则映射 \(\mathsf{pr}_1\) 诱导出从 \(\Gamma\) 到 X 的同构，并且 \(X\)\(\mathrm{T}_{(x,y)}(\Gamma)\) 与 \(\mathrm{T}_x(f)\) 的图等同。

特别地，\(X \times X\) 的对角线是 \(X \times X\) 的子流形。

5.8.10. 设 Y 是一个流形，\((f_{i})_{i\in I}\) 是 Y 上的有限个态射函数组成的族。令 X 为满足对每个 i 都有 \(x\in Y\) 的 \(f_{i}(x)=0\) 的集合。假设下列条件成立：

\begin{enumerate}
\def\labelenumi{(\Alph{enumi})}
\setcounter{enumi}{9}
\tightlist
\item
  对每个 \(x \in X\)，微分 \(d_x f_i\) 在 \(T_x'(Y)\) 中线性无关。
\end{enumerate}

则 X 是 Y 的闭子流形，切空间 \(T_{x}(X)\) 是由满足对 I 中每个 i 都有 \(T_{x}(Y)\) 的 \(\alpha\) 所组成的 \(\alpha.f_{i}=0\) 的子空间。此外，X 在 Y 中每一点处的余维都等于 Card (I)。

5.8.11.（理想的简单点）。设 \(\mathfrak{a}\) 是多项式代数 \(\mathrm{K}[\mathrm{X}_1,\ldots ,\mathrm{X}_n]\) 的一个理想。若  中的点 \(x = (x_{1},\dots,x_{n})\) 对所有 \(\mathbf{K}^n\) 都满足 f(x) = 0，则称其为 \(\mathfrak{a}\) 的零点。若 \(f(x) = 0\)\(f\in \mathfrak{a}\)\(x\in \mathbf{K}^n\)，以 \(\mathbf{S}_x\) 表示由分式  构成的 \(\mathrm{K}(\mathrm{X}_1,\dots,\mathrm{X}_n)\) 的子代数，其中 \(f / g\)\(f,g\in \mathrm{K}[\mathrm{X}_1,\dots,\mathrm{X}_n]\) 且 \(g(x)\neq 0\)；以 \(\mathfrak{a}_x\) 表示 \(\mathbf{S}_x\) 中由 \(\mathfrak{a}\) 生成的理想。若点 \(\mathbf{S}_x\)\(x\) 是 \(\mathfrak{a}\) 的零点且满足下列条件，则称 x 为 \(\mathfrak{a}\) 的简单零点：

\begin{enumerate}
\def\labelenumi{(\Alph{enumi})}
\setcounter{enumi}{18}
\tightlist
\item
  存在一个由 \((f_1, \ldots, f_m)\) 中元素组成的有限序列 \(\mathfrak{a}\)，它们生成理想 \(\mathfrak{a}_x\)，且它们在 \(x\) 处的微分线性无关。（该条件等价于局部环 \(\mathbf{S}_x / \mathfrak{a}_x\) 是正则的（交换代数，即将出版）。）
\end{enumerate}

令 \(Z\)（分别为 \(Z_s\)）为 \(\mathfrak{a}\) 的零点集（分别为简单零点集）。集合 \(Z\) 在 \(K^n\) 中是闭的，\(Z_s\) 在 \(Z\) 中是开的，而 \(Z_s\) 是 \(K^{n}\) 的子流形。若 \(x \in Z_{s}\)，则理想 \(\mathrm{K}[\mathrm{X}_{1}, \ldots, \mathrm{X}_{n}] \cap \mathfrak{a}_{x}\) 由在 \(Z_{s}\) 中 x 的某个邻域上消失的多项式 f 构成。

令 \(\overline{\mathfrak{a}}\) 为在 Z 上消失的多项式构成的理想。若 K 是代数闭域，则 \(\mathfrak{a}\) 的简单零点集在 Z 中稠密。

5.8.12. 设 X 是一个流形，L 是由有序对 \((x, Z)\) 组成的集合，其中 x 是 X 的一点，Z 是包含 x 的 X 的子流形。给定两个有序对 \(\pi = (x, Z)\) 和 \(\pi' = (x', Z')\)，用 \(R\{\pi, \pi'\}\) 表示如下关系：

“我们有 \(x = x'\)，且存在 \(x\) 的一个邻域 U，使得 \(\mathbf{U} \cap \mathbf{Z} = \mathbf{U} \cap \mathbf{Z'}\)。于是 R 是 L 上的等价关系；把有序对 \(\gamma_x(Z)\) 的等价类记为 \((x, Z) \in L\)。在集合 \(J = L/R\) 上存在唯一的流形结构，满足下列条件：

对于 X 的每个子流形 Z，映射 \(x \mapsto \gamma_{x}(Z)\) 从 Z 到 J 是从 Z 到 J 的一个开子流形上的同构。

称 J 为 X 的子流形芽的流形（参见~一般拓扑，第~I~章，第 4 \(^{e}\) 版，§ 6，第 10 号）。

由  定义的映射 \(\rho: J \to X\) 是一个浸入；称其为 J 到 X 的典范浸入。\(\rho(\gamma_x(Z)) = x\)\(J\)\(X\)

若 X 是有限维的分离解析流形，则 J 也如此。

